% !TEX root = ../../../main.tex
% 所属：第一部分：单变量分析学 / 第4章：数列极限的应用与一些反例
% 节：4.4 $\pi,e,\gamma$的定义，性质及应用
% 正文按原文件顺序保存；由 ../chapter4.tex 调用。
\section{$\pi,e,\gamma$的定义，性质及应用}
\begin{proposition}
	数列$\left\{n\sin\dfrac{180^\circ}{n}\right\}$收敛.
\end{proposition}
\begin{proof}
	
	令$L_n=n\sin\dfrac{180^\circ}{n}$,设 $t = \dfrac{180^\circ}{n(n+1)}$，对于 $n \geqslant 3$，且 $nt \leqslant 45^\circ$，于是
	
	\[
	\tan nt = \frac{\tan(n-1)t + \tan t}{1 - \tan(n-1)t \tan t}\geqslant\tan(n-1)t + \tan t\geqslant\cdots
	\geqslant n\tan t
	\]
	
	从而，
	\[\sin(n+1)t = \sin nt \cos t + \sin t \cos nt = \sin nt \cos t \left(1 + \frac{\tan t}{\tan nt}\right) \leqslant \frac{n+1}{n} \sin nt\]
	因此对于 $n \geqslant 3$，有
	\[
	L_n = n \sin \frac{180^\circ}{n} \leqslant (n+1) \sin \frac{180^\circ}{n+1} = L_{n+1}
	\]
	单位圆内接正 $n$ 边形的面积
	
	\[
	S_n = n \sin \frac{180^\circ}{n} \cos \frac{180^\circ}{n} < 4
	\]
	
	对于 $n \geqslant 3$，有
	
	\[
	L_n = n \sin \frac{180^\circ}{n} < \frac{4}{\cos 60^\circ} = 8
	\]
	
	所以，$\{L_n\}$ 单调递增且有上界，因而收敛.
	
\end{proof}
\begin{definition}
	数列$\left\{n\sin\dfrac{180^\circ}{n}\right\}$的收敛值定义为$\pi$
\end{definition}
使用python,输入以下代码
\begin{verbatim}
	"""
	用数列 {n * sin(180°/n)} 的收敛值计算 π
	"""
	
	import math
	def calculate_pi(n: int) -> float:
	return n * math.sin(math.pi / n)
	def main():
	print("=" * 65)
	print("用数列 {n * sin(180°/n)} 的收敛值计算 π")
	print("=" * 65)
	print(f"{'n':<12} {'近似值':<20} {'与 π 的误差':<20}")
	print("-" * 65)
	n_values = [1, 2, 3, 5, 10, 50, 100, 500, 1000, 5000, 10000, 50000, 100000, 1000000]
	for n in n_values:
	approx_pi = calculate_pi(n)
	error = abs(approx_pi - math.pi)
	print(f"{n:<12} {approx_pi:<20.15f} {error:<20.15e}")
	print("\n" + "=" * 65)
	print(f"π 的真实值: {math.pi:.15f}")
	print("=" * 65)
	if __name__ == "__main__":
	main()
\end{verbatim}
能够得到以下结果
\begin{verbatim}
================================================================
用数列 {n * sin(180°/n)} 的收敛值计算 pi
=================================================================
n            近似值                  与 pi 的误差             
-----------------------------------------------------------------
1            0.000000000000000    3.141592653589793e+00
2            2.000000000000000    1.141592653589793e+00
3            2.598076211353316    5.435164422364771e-01
5            2.938926261462366    2.026663921274272e-01
10           3.090169943749474    5.142270984031905e-02
50           3.139525976465669    2.066677124124450e-03
100          3.141075907812829    5.167457769639228e-04
500          3.141571982779475    2.067081031764673e-05
1000         3.141587485879564    5.167710229514455e-06
5000         3.141592446881286    2.067085071999486e-07
10000        3.141592601912665    5.167712791021017e-08
50000        3.141592651522708    2.067085169699112e-09
100000       3.141592653073022    5.167715144693830e-10
1000000      3.141592653584626    5.167422045815329e-12

=================================================================
pi的真实值: 3.141592653589793
=================================================================	
\end{verbatim}
这样我们可以使用$\pi$的近似值$3.14159265358979323846...$

关于$\pi$是有理数还是无理数这一点，我们之后再说.

并且我们可以定义出弧度制,
$$\text{半个圆周对应的圆周角}:=\pi$$
在弧度制下$$\lim\limits_{n\to\infty}\frac{\sin(\pi/n)}{\pi/n}=1$$


\begin{proposition}
	设 $x_n = \left(1 + \dfrac{1}{n}\right)^n$，$y_n = \left(1 + \dfrac{1}{n}\right)^{n+1}$，证明 $\{x_n\}$ 单调递增，$\{y_n\}$ 单调递减,并且收敛于同一极限.
\end{proposition}
\begin{proof}
	对于 $\{x_n\}$，
	\[
	x_n = \left(1 + \frac{1}{n}\right)^n\cdot 1 < \left(\dfrac{n(1 + \dfrac{1}{n})+1}{n+1}\right)^{n+1} = \left(1 + \dfrac{1}{n+1}\right)^{n+1} < x_{n+1}
	\]
	
	对于 $\{y_n\}$，
	\[
	\frac{1}{y_n} = \left(\dfrac{n}{n+1}\right)^{n+1}\cdot 1 < \left(\dfrac{(n+1)\dfrac{n}{n+1}+1}{n+2}\right)^{n+2} = \left(\dfrac{n+1}{n+2}\right)^{n+2} < \frac{1}{y_{n+1}}
	\]
	
	因此，$\{x_n\}$ 单调递增，$\{y_n\}$ 单调递减.又因为 $x_n = y_n(1+\dfrac{1}{n})$，故 $\{x_n\}$，$\{y_n\}$ 收敛于同一极限.
	
	
\end{proof}
\begin{definition}
	\[
	\lim_{n \to \infty} \left(1 + \frac{1}{n}\right)^n = \lim_{n \to \infty} \left(1 + \frac{1}{n}\right)^{n+1} = e
	\]
	
	称为自然对数的底数,自然对数记为
	\[
	\ln x = \log_e x
	\]
	
\end{definition}
使用以下代码
\begin{verbatim}
	import math
	def calculate_e_lower(n: int) -> float:
	return (1 + 1/n) ** n
	def calculate_e_upper(n: int) -> float:
	return (1 + 1/n) ** (n + 1)
	def main():
	print("=" * 80)
	print("计算 e")
	print("=" * 80)
	print(f"{'n':<10} {'(1+1/n)^n':<20} {'(1+1/n)^(n+1)':<20} 
	{'误差(下)':<18} {'误差(上)':<18}")
	print("-" * 80)
	n_values = [1, 2, 5, 10, 50, 100, 500, 1000, 5000, 10000, 50000, 100000, 1000000]
	for n in n_values:
	lower = calculate_e_lower(n)
	upper = calculate_e_upper(n)
	err_lower = abs(lower - math.e)
	err_upper = abs(upper - math.e)
	print(f"{n:<10} {lower:<20.15f} {upper:<20.15f} {err_lower:<18.12e} {err_upper:<18.12e}")
	print("\n" + "=" * 80)
	print(f"e 的真实值: {math.e:.15f}")
	print("=" * 80)
	# 验证不等式关系
	print("\n验证不等式关系 (1+1/n)^n < e < (1+1/n)^(n+1)：")
	for n in [1, 2, 10, 100, 1000]:
	lower = calculate_e_lower(n)
	upper = calculate_e_upper(n)
	print(f"  n={n:6d}:  {lower:.15f}  <  {math.e:.15f}  <  {upper:.15f}  ✓")
	if __name__ == "__main__":
	main()
\end{verbatim}
输出结果如下
\begin{verbatim}
	================================================================================
	计算 e
	================================================================================
	n          (1+1/n)^n            (1+1/n)^(n+1)        误差(下)              误差(上)             
	--------------------------------------------------------------------------------
	1          2.000000000000000    4.000000000000000    7.182818284590e-01 1.281718171541e+00
	2          2.250000000000000    3.375000000000000    4.682818284590e-01 6.567181715410e-01
	5          2.488319999999999    2.985983999999999    2.299618284590e-01 2.677021715410e-01
	10         2.593742460100002    2.853116706110002    1.245393683590e-01 1.348348776510e-01
	50         2.691588029073608    2.745419789655080    2.669379938544e-02 2.713796119603e-02
	100        2.704813829421528    2.731861967715744    1.346799903752e-02 1.358013925670e-02
	500        2.715568520651728    2.720999657693032    2.713307807317e-03 2.717829233987e-03
	1000       2.716923932235594    2.719640856167829    1.357896223452e-03 1.359027708784e-03
	5000       2.718010050101555    2.718553652111575    2.717783574901e-04 2.718236525299e-04
	10000      2.718145926824926    2.718417741417608    1.359016341196e-04 1.359129585627e-04
	50000      2.718254646126732    2.718309011219655    2.718233231302e-05 2.718276060953e-05
	100000     2.718268237192297    2.718295419874670    1.359126674760e-05 1.359141562451e-05
	1000000    2.718280469095753    2.718283187376222    1.359363291709e-06 1.358917177008e-06
	
	================================================================================
	e 的真实值: 2.718281828459045
	================================================================================
	
	验证不等式关系 (1+1/n)^n < e < (1+1/n)^(n+1)：
	n=     1:  2.000000000000000  <  2.718281828459045  <  4.000000000000000  ✓
	n=     2:  2.250000000000000  <  2.718281828459045  <  3.375000000000000  ✓
	n=    10:  2.593742460100002  <  2.718281828459045  <  2.853116706110002  ✓
	n=   100:  2.704813829421528  <  2.718281828459045  <  2.731861967715744  ✓
	n=  1000:  2.716923932235594  <  2.718281828459045  <  2.719640856167829  ✓
\end{verbatim}
\begin{proposition}
	设$S_n=\sum\limits_{k=0}^{n}\dfrac{1}{k!}$，那么$e=\lim\limits_{n \to \infty}S_n$
\end{proposition}
\begin{proof}
	对于 $\{x_n\}$，由二项式定理，
	\[
	\left(1 + \frac{1}{n}\right)^n = \sum_{k=0}^n C_n^k \frac{1}{n^k}
	\]
	
	对于任意 $k \leqslant n$，都有
	\[
	C_n^k = \frac{n(n-1)\cdots(n-k+1)}{k!} \leqslant \frac{n^k}{k!}
	\]
	
	$\forall n \in\Z_{+}$，有
	\[
	\left(1 + \frac{1}{n}\right)^n = \sum_{k=0}^n C_n^k \frac{1}{n^k} < \sum_{k=0}^n \frac{1}{k!}
	\]
	
	对于 $\{y_n\}$，由二项式公式，
	\[
	\left(1 + \frac{1}{n}\right)^{n+1} = \sum_{k=0}^{n+1} C_{n+1}^k \frac{1}{n^k}
	\]
	
	对于任意 $k \leqslant n$，都有
	\[
	C_{n+1}^k = \frac{(n+1)n(n-1)\cdots(n+1-k+1)}{k!} = \frac{n^k}{k!} \prod_{j=0}^{k-1} \left(1 + \frac{1 - j}{n}\right) \geqslant \frac{n^k}{k!}
	\]
	
	得对所有 $n$，有
	\[
	\sum_{k=0}^n \frac{1}{k!} \leqslant \left(1 + \frac{1}{n}\right)^{n+1} = \sum_{k=0}^{n+1} C_{n+1}^k \frac{1}{n^k}
	\]
	
	由夹逼定理
	\[
	\lim_{n \to \infty} \sum_{k=0}^n \frac{1}{k!} = \sum_{k=0}^\infty \frac{1}{k!} =\lim\limits_{n \to \infty}(1+\frac{1}{n})^n= e
	\]
\end{proof}
这样我们给出了$e$的级数定义.这里我们把级数当成一种特殊的数列来考虑，关于级数的理论我们放在下一章讨论.
\begin{theorem}
	$e$是一个无理数.
\end{theorem}
\begin{proof}
	假设 $e$ 有理，记 $e=\dfrac{p}{q}$（最简分数，$p, q \in \mathbb{Z}^{+}$）.取 $q$ 固定，考虑部分和
	
	$$
	S_q=\sum_{k=0}^q \frac{1}{k!} .
	$$
	
	
	注意对 $k \leqslant q$ 有 $q!/ k!$ 为整数，所以
	
	$$
	q!S_q=\sum_{k=0}^q \frac{q!}{k!} \in \mathbb{Z}
	$$
	
	
	另一方面用 $e$ 的级数表示
	
	$$
	q!e=q!S_q+q!\sum_{k=q+1}^{\infty} \frac{1}{k!}
	$$
	
	
	左边 $q!e=\dfrac{q!p}{q}=(q-1)!p$ 是整数，由此可知尾和
	
	$$
	T:=q!\sum_{k=q+1}^{\infty} \frac{1}{k!}
	$$
	
	
	也是整数.我们来估计这个 $T$ ：
	对 $m \geqslant 1$ ，
	
	$$
	\frac{q!}{(q+m)!}=\frac{1}{(q+1)(q+2) \cdots(q+m)} \leqslant \frac{1}{(q+1)^m}
	$$
	
	
	因此
	
	$$
	T=\sum_{m=1}^{\infty} \frac{q!}{(q+m)!} \leqslant \sum_{m=1}^{\infty} \frac{1}{(q+1)^m}=\frac{1}{q}
	$$
	
	
	又显然每一项均为正，所以 $T>0$ .于是对于 $q \geqslant 2$（若 $q=1$ 则 $e$ 若是整数显然不成立），我们得到
	
	$$
	0<T<\frac{1}{q} \leqslant 1
	$$
	
	
	即 $T$ 是一个严格在 0 与 1 之间的正整数——这是不可能的.矛盾.
	因此假设错误，$e$ 不能是有理数，亦即 $e$ 是无理数
	
\end{proof}

\begin{corollary}
	$$\sum\limits_{k=0}^{n}\dfrac{1}{k!}+\frac{1}{(n+1)!}<e<\sum\limits_{k=0}^{n}\dfrac{1}{k!}+\frac{1}{n\cdot n!}$$
\end{corollary}
\begin{proof}
	$$\sum\limits_{k=0}^{n}\dfrac{1}{k!}+\frac{1}{(n+1)!}<e$$由$e$的级数定义是显然的
	
	
	$$\begin{aligned}
		e - \sum_{k=0}^n \frac{1}{k!} = \sum_{k=n+1}^\infty \frac{1}{k!} &= \frac{1}{(n+1)!} \left(1 + \frac{1}{n+2} + \frac{1}{(n+2)(n+3)} + \cdots \right)\\
		&< \frac{1}{(n+1)!} \left(1 + \frac{1}{n+1} + \frac{1}{(n+1)^2} + \cdots \right)\\
		&= \frac{1}{(n+1)!} \cdot \frac{n+1}{n} = \frac{1}{n \cdot n!}
	\end{aligned}$$
	
\end{proof}
\begin{example}
	证明:
	
	(1) $\lim\limits_{n \to \infty} \left(1 + \dfrac{k}{n}\right)^n = e^k$
	
	(2) $\lim\limits_{n \to \infty} \left(1 + \dfrac{1}{n + k}\right)^n = e$
	
	(3) $\lim\limits_{n \to \infty} \left(1 + \dfrac{1}{n^2}\right)^n = 1$
	
	(4) $\lim\limits_{n \to \infty} \left(1 + \dfrac{1}{n} - \dfrac{1}{n^2}\right)^n = e$
	
	
\end{example}
\begin{proof}
	
	(1) $\lim\limits_{n \to \infty} \left[\left(1 + \dfrac{1}{n/k}\right)^{n/k}\right]^k = e^k$
	
	(2) $\lim\limits_{n \to \infty} \left(1 + \dfrac{1}{n + k}\right)^{n + k} \cdot \left(1 + \dfrac{1}{n + k}\right)^{-k} = e$
	
	(3) $\lim\limits_{n \to \infty} \left[\left(1 + \dfrac{1}{n^2}\right)^{n^2}\right]^{1/n} = \lim\limits_{n \to \infty} e^{1/n} = 1$
	
	(4) 对于 $n \geqslant 2$，有
	\[
	\left(1 + \frac{1}{n+2}\right)^n \leqslant \left(1 + \frac{1}{n} - \frac{1}{n^2}\right)^n < \left(1 + \frac{1}{n}\right)^n
	\]
	
	由夹逼定理，
	\[
	\lim_{n \to \infty} \left(1 + \frac{1}{n} - \frac{1}{n^2}\right)^n = e
	\]
	
\end{proof}
\begin{proposition}\label{调和级数}
	设 $a_n = 1 + \dfrac{1}{2^p} + \dfrac{1}{3^p} + \cdots + \dfrac{1}{n^p}$，当 $p > 1$，$\{a_n\}$ 收敛，当$p \leqslant 1$，$\{a_n\}$ 发散.
	
	
\end{proposition}
\begin{proof}
	显然的，$\{a_n\}$ 单调递增且有上界.
	
	对于 $p > 1$，
	$$\begin{aligned}
		&\frac{1}{2^p} + \frac{1}{3^p} < \frac{1}{2^{p}} + \frac{1}{2^{p}} = \frac{1}{2^{p-1}} = r\\
		&\frac{1}{4^p} + \frac{1}{5^p} + \cdots + \frac{1}{7^p} < 4 \cdot \frac{1}{4^p} = \frac{1}{4^{p-1}} = r^2\\
		&\cdots\\
		&\frac{1}{2^{kp}} + \cdots + \frac{1}{(2^{k+1}-1)^p} < 2^k \cdot \frac{1}{2^{kp}} = \frac{1}{2^{kp-k}} = r^k
	\end{aligned}$$
	
	
	因此，
	\[
	a_n \leqslant a_{2^n-1} < 1 + r + r^2 + \cdots + r^{n-1} < \frac{1}{1-r}
	\]
	
	所以 $p > 1$ 时，$\{a_n\}$ 有界，则 $\{a_n\}$ 收敛.
	
	对于 $p \leqslant 1$，
	$$\begin{aligned}
		&\frac{1}{2^p}  \geqslant \frac{1}{2} \\
		&\frac{1}{3^p}+\frac{1}{4^p}  > 2\cdot\frac{1}{4}\\
		&\cdots\\
		&\frac{1}{(2^{k}+1)^{p}} + \cdots + \frac{1}{(2^{k+1})^{p}} > \frac{2^{k}}{2^{k+1}} = \frac{1}{2}
	\end{aligned}$$
	
	这样$$a_{2^n}>1+\frac{n}{2}$$
	
	因此，$\{a_{2^n}\}$ 是正无穷大量，故$\{a_n\}$ 发散
	
	
\end{proof}
特别地，$n=1$时$a_n=\sum\limits_{k=1}^n\dfrac{1}{k}$发散，称$\sum\limits_{k=1}^{\infty}\dfrac{1}{k}$为调和级数，同时我们定义$$\zeta(n)=\sum\limits_{k=1}^{\infty}\dfrac{1}{k^n}$$这是经典的黎曼Zeta函数，我们在复变函数的部分研究他的性质，它在解析数论中有重要的作用.
\begin{proposition}
	设 $b_n = 1 + \dfrac{1}{2} + \dfrac{1}{3} + \cdots + \dfrac{1}{n} - \ln n$，则 $\{b_n\}$ 是收敛数列.
	\label{defofeucon}
\end{proposition}
\begin{proof}
	由不等式
	\[
	\left(1 + \frac{1}{n}\right)^n < e < \left(1 + \frac{1}{n}\right)^{n+1}
	\]
	
	取自然对数，
	\[
	n \ln \left(1 + \frac{1}{n}\right) < 1 < (n+1) \ln \left(1 + \frac{1}{n}\right)
	\]
	
	即
	\[
	\frac{1}{n+1} < \ln \frac{n+1}{n} < \frac{1}{n}
	\]
	
	考虑差分：
	\[
	b_{n+1} - b_n = \frac{1}{n+1} - \ln(n+1) + \ln n < 0
	\]
	
	因此，
	\[
	b_n = 1 + \frac{1}{2} + \frac{1}{3} + \cdots + \frac{1}{n} - \ln n >\ln2-\ln1+\ln3-\ln2+\cdots+\ln(n+1)-\ln n-\ln n= \ln \frac{n+1}{n} > 0
	\]
	
	$\{b_n\}$ 单调递减且有下界，从而收敛.
\end{proof}
\begin{definition}
	将$b_n = 1 + \dfrac{1}{2} + \dfrac{1}{3} + \cdots + \dfrac{1}{n} - \ln n$，则 $\{b_n\}$的收敛值定义为$\gamma$，称为欧拉-马歇罗尼常数
\end{definition}
使用以下代码
\begin{verbatim}
	import math
	def harmonic_partial_sum(n: int) -> float:
	total = 0.0
	for k in range(1, n + 1):
	total += 1.0 / k
	return total
	def calculate_gamma(n: int) -> float:
	return harmonic_partial_sum(n) - math.log(n)
	def main():
	gamma_true = 0.57721566490153286060651209008240243104215933593992
	print("=" * 65)
	print("用数列 b_n = H_n - ln(n) 计算欧拉-马歇罗尼常数 gamma")
	print("=" * 65)
	print(f"{'n':<12} {'b_n':<22} {'与 gamma 的误差':<20}")
	print("-" * 65)
	n_values = [10, 100, 500, 1000, 5000, 10000, 50000, 100000, 500000, 1000000]
	for n in n_values:
	gamma_approx = calculate_gamma(n)
	error = abs(gamma_approx - gamma_true)
	print(f"{n:<12} {gamma_approx:<22.15f} {error:<20.12e}")
	print("\n" + "=" * 65)
	print(f"gamma 的高精度参考值: {gamma_true:.15f}")
	print("=" * 65)
	if __name__ == "__main__":
	main()
	
\end{verbatim}
得到
\begin{verbatim}
	=================================================================
	用数列 b_n = H_n - ln(n) 计算欧拉-马歇罗尼常数 gamma
	=================================================================
	n            b_n                    与 gamma 的误差         
	-----------------------------------------------------------------
	10           0.626383160974208      4.916749607268e-02  
	100          0.582207331651529      4.991666749996e-03  
	500          0.578215331568328      9.996666667956e-04  
	1000         0.577715581568206      4.999166666736e-04  
	5000         0.577315661568166      9.999666663318e-05  
	10000        0.577265664068165      4.999916663173e-05  
	50000        0.577225664868221      9.999966688379e-06  
	100000       0.577220664893106      4.999991573551e-06  
	500000       0.577216664901298      9.999997655497e-07  
	1000000      0.577216164900715      4.999991823906e-07  
	
	=================================================================
	gamma 的高精度参考值: 0.577215664901533
	=================================================================
\end{verbatim}
这样我们得到$\gamma$的近似值$0.57721566490...$
\begin{corollary}
	(1) $\lim\limits_{n \to \infty} \left(\dfrac{1}{n+1} + \cdots + \dfrac{1}{2n}\right) = \ln 2$
	
	(2) $\lim\limits_{n \to \infty} \left[1 - \dfrac{1}{2} + \dfrac{1}{3} - \cdots + (-1)^{n+1} \dfrac{1}{n}\right] = \ln 2$
\end{corollary}
\begin{proof}
	(1) 设 $C_n = \dfrac{1}{n+1} + \cdots + \dfrac{1}{2n}$，由\ref{defofeucon}
	\[
	C_n = b_{2n} - b_n + \ln(2n) - \ln n = b_{2n} - b_n + \ln 2
	\]
	
	又$\lim\limits_{n \to \infty} b_n = \gamma$，则 $\lim\limits_{n \to \infty} b_{2n} = \gamma$，所以
	\[
	\lim_{n \to \infty} C_n = \lim_{n \to \infty} \left(\frac{1}{n+1} + \cdots + \frac{1}{2n}\right) = \ln 2
	\]
	
	(2) 设 $d_n = 1 - \dfrac{1}{2} + \dfrac{1}{3} - \cdots + \dfrac{1}{2n-1} - \dfrac{1}{2n}$，
	$$\begin{aligned}
		&	b_{2n} = 1 + \frac{1}{2} + \frac{1}{3} + \cdots + \frac{1}{2n} - \ln(2n)\\
		&	b_n = 1 + \frac{1}{2} + \frac{1}{3} + \cdots + \frac{1}{n} - \ln n\\
		&b_{2n} - b_n = 1 - \frac{1}{2} + \frac{1}{3} - \cdots - \frac{1}{2n} - \ln(2n) + \ln n\\
		&d_{2n} = b_{2n} - b_n + \ln(2n) - \ln n = b_{2n} - b_n + \ln 2 \to \ln 2\\
		&d_{2n+1} = d_{2n} + \frac{1}{2n+1} \to \ln 2
	\end{aligned}$$
	
	因此，$\lim\limits_{n \to \infty} d_n = \ln 2$.
	
\end{proof}
